%% =======================================================
%% Actividad 15: ¿Qué significa contrastar una hipótesis?
%%
%% Compilar desde actividades/: ./compilar.sh act15
%% Esta actividad no tiene código: el mapa conceptual está dibujado con TikZ.
%% =======================================================

%% ---- CONFIGURACIÓN ------------------------

\documentclass[11pt,letterpaper]{article}

\usepackage[indice]{tecnm}
\usepackage[T1]{fontenc}
\usepackage{amsmath}
\usepackage{xcolor}
\usepackage{tikz}
\usetikzlibrary{arrows.meta,positioning,calc}
\setlength{\emergencystretch}{3em} % Evita líneas desbordadas por palabras en \texttt

%% ---- DATOS DE LA MATERIA ------------------------

\materia{Modelos Estadísticos para el Aprendizaje \\ Automático y Ciencia de Datos}
\actividad{Tema 3.1 Conceptos de las Pruebas de Hipótesis \\ Actividad 15 - ¿Qué significa contrastar una hipótesis?}

\profesor{Dr. José Ismael Ojeda Campaña}
%\fecha{15 de octubre de 2026}    % opcional; por defecto es la fecha de hoy

\begin{document}
\portada
\indicegeneral

%% ---- CONTENIDO ---------------------------------

\section{Introducción}

Una muestra siempre contiene variabilidad, así que una diferencia observada puede deberse a un efecto real o solo al azar. Las pruebas de hipótesis son el procedimiento estadístico que permite decidir entre ambas explicaciones: evalúan si una afirmación sobre un parámetro de la población es razonable a la luz de los datos de una muestra \parencite{tecnm2026_documento_fuente}. Por eso son la base para tomar decisiones con datos, por ejemplo al comparar dos modelos de aprendizaje automático.

En esta actividad se elabora un mapa conceptual que relaciona los cinco conceptos centrales del contraste de hipótesis: la hipótesis nula, la hipótesis alternativa, el nivel de significancia, el valor $p$ y los errores tipo I y tipo II. Primero se definen los conceptos, después se presenta el mapa con sus conexiones y por último se explican esas conexiones con ejemplos básicos.

\section{Marco conceptual}

\subsection{Hipótesis nula y alternativa}

La hipótesis nula, $H_0$, se formula como la hipótesis en la que no hay diferencia o no hay efecto. La hipótesis alternativa, $H_1$, es cualquier hipótesis que difiera de $H_0$ y, por lo regular, representa lo que se quiere demostrar \parencite{tecnm2026_documento_fuente}. Ambas son proposiciones opuestas sobre el mismo parámetro: si una es cierta, la otra es falsa. Por ejemplo, al decidir si un procedimiento es mejor que otro, $H_0$ afirma que no existe diferencia entre ellos y $H_1$ afirma que no son iguales.

\subsection{Nivel de significancia}

El nivel de significancia, $\alpha$, es la probabilidad máxima de cometer un error tipo I en la prueba. Sus valores más comunes son 0.05 y 0.01. Con $\alpha = 0.05$, aproximadamente en 5 de cada 100 ocasiones se rechazaría la hipótesis nula cuando debería ser aceptada \parencite{tecnm2026_documento_fuente,spiegel2014}. El valor de $\alpha$ lo fija quien diseña la prueba antes de ver los datos, y funciona como el umbral de decisión.

\subsection{Valor $p$}

El valor $p$ es un número que indica qué tan compatibles son los datos con la hipótesis nula \parencite{tecnm2026_documento_fuente}. Formalmente, es la probabilidad de obtener un resultado tan extremo como el observado, o más, suponiendo que $H_0$ es verdadera. Un valor $p$ pequeño significa que los datos serían poco esperables si $H_0$ fuera cierta. No es la probabilidad de que $H_0$ sea verdadera ni mide el tamaño del efecto.

\subsection{Errores tipo I y tipo II}

Si se rechaza una hipótesis cuando debería ser aceptada, se comete un error tipo I; si se acepta una hipótesis que debería ser rechazada, se comete un error tipo II \parencite{tecnm2026_documento_fuente}. La probabilidad del error tipo I es $\alpha$ y la del tipo II se denota $\beta$. Para un tamaño de muestra dado, disminuir un tipo de error aumenta el otro, y aumentar el tamaño de la muestra es la única forma de reducir ambos al mismo tiempo, aunque no siempre sea posible \parencite{tecnm2026_documento_fuente}.

\section{Descripción de la actividad}

\textbf{Objetivo:} Identificar los fundamentos conceptuales de las pruebas de hipótesis y su importancia en la toma de decisiones basadas en datos \parencite{tecnm2026_documento_fuente}.

Se leyó el texto introductorio sobre el proceso de inferencia estadística y se elaboró un mapa conceptual que muestra la relación entre la hipótesis nula ($H_0$), la hipótesis alternativa ($H_1$), el nivel de significancia ($\alpha$), el valor $p$ y los errores tipo I y tipo II, con conexiones lógicas y ejemplos básicos. El producto es este documento en PDF.

\section{Desarrollo}

\subsection{Mapa conceptual}

La Figura~\ref{fig:mapa} muestra el mapa. Se lee de arriba hacia abajo: primero se plantean las hipótesis, luego se compara la evidencia (valor $p$) con el umbral ($\alpha$), y la decisión resultante puede ser correcta o puede ser uno de los dos errores. Las flechas llevan la palabra de enlace que explica cada relación.

\begin{figure}[htbp]
\centering
\begin{tikzpicture}[
    font=\footnotesize,
    >={Stealth[length=2.2mm]},
    concepto/.style={draw=tecnmazul, rounded corners=3pt, thick, align=center,
      inner sep=4pt, text width=3.6cm, fill=tecnmazul!8},
    central/.style={concepto, fill=tecnmazul, text=white, text width=4.4cm},
    decision/.style={concepto, fill=gray!15, draw=gray!70!black},
    error/.style={concepto, fill=red!8, draw=red!60!black, text width=4.4cm},
    flecha/.style={->, thick, draw=gray!70!black},
    enlace/.style={font=\scriptsize\itshape, fill=white, inner sep=1.5pt, align=center}
  ]
  % Nivel 0
  \node[central] (prueba) at (0,0) {\textbf{Prueba de hipótesis}\\ decide con datos de una muestra};

  % Nivel 1: hipótesis
  \node[concepto] (h0) at (-4.3,-2.6) {\textbf{Hipótesis nula $H_0$}\\ no hay diferencia ni efecto};
  \node[concepto] (h1) at (4.3,-2.6) {\textbf{Hipótesis alternativa $H_1$}\\ hay diferencia o efecto};

  % Nivel 2: evidencia y umbral
  \node[concepto] (p) at (-4.3,-5.5) {\textbf{Valor $p$}\\ compatibilidad de los datos con $H_0$};
  \node[concepto] (alfa) at (4.3,-5.5) {\textbf{Nivel de significancia $\alpha$}\\ umbral; probabilidad máxima de error tipo I};

  % Nivel 3: regla de decisión
  \node[decision, text width=4.6cm] (regla) at (0,-8.3) {\textbf{Regla de decisión}\\ comparar $p$ con $\alpha$};

  % Nivel 4: decisión
  \node[decision] (nr) at (-4.3,-11.1) {\textbf{No rechazar $H_0$}\\ los datos no bastan para contradecirla};
  \node[decision] (r) at (4.3,-11.1) {\textbf{Rechazar $H_0$}\\ los datos apoyan $H_1$};

  % Nivel 5: errores
  \node[error] (e2) at (-4.3,-14.2) {\textbf{Error tipo II}\\ no rechazar $H_0$ siendo falsa\\ probabilidad $\beta$};
  \node[error] (e1) at (4.3,-14.2) {\textbf{Error tipo I}\\ rechazar $H_0$ siendo verdadera\\ probabilidad $\alpha$};

  % Enlaces
  \draw[flecha] (prueba.south) -- ++(0,-0.7) -| (h0.north) node[enlace, pos=0.72] {parte de};
  \draw[flecha] (prueba.south) -- ++(0,-0.7) -| (h1.north) node[enlace, pos=0.72] {parte de};
  \draw[<->, thick, draw=gray!70!black] (h0) -- (h1) node[enlace, midway] {son\\opuestas};
  \draw[flecha] (h0) -- (p) node[enlace, midway] {se calcula\\suponiéndola cierta};
  \draw[flecha] (p.south) -- ++(0,-0.5) -| (regla.north west) node[enlace, pos=0.3] {evidencia};
  \draw[flecha] (alfa.south) -- ++(0,-0.5) -| (regla.north east) node[enlace, pos=0.3] {umbral};
  \draw[flecha] (regla.south) -- ++(0,-0.4) -| (nr.north) node[enlace, pos=0.25] {si $p > \alpha$};
  \draw[flecha] (regla.south) -- ++(0,-0.4) -| (r.north) node[enlace, pos=0.25] {si $p \le \alpha$};
  \draw[flecha] (r) -- (e1) node[enlace, midway] {si $H_0$ era\\verdadera};
  \draw[flecha] (nr) -- (e2) node[enlace, midway] {si $H_0$ era\\falsa};
  \draw[<->, thick, dashed, draw=red!60!black] (e2) -- (e1)
    node[enlace, midway] {con $n$ fijo, bajar uno\\sube el otro};
  \draw[flecha] (alfa.east) -- ++(0.75,0) |- (e1.east) node[enlace, pos=0.2, text width=1.5cm] {fija el\\error tipo I};
\end{tikzpicture}
\caption{Mapa conceptual de las pruebas de hipótesis.}
\label{fig:mapa}
\end{figure}

\subsection{Descripción de los conceptos con ejemplos básicos}

La Tabla~\ref{tab:conceptos} resume cada concepto con un ejemplo. Los ejemplos son ilustrativos y usan la comparación de un modelo nuevo contra el modelo actual, que es el caso típico en aprendizaje automático. Los valores numéricos son hipotéticos.

\begin{table}[htbp]
\centering
\caption{Conceptos de la prueba de hipótesis y un ejemplo básico de cada uno.}
\label{tab:conceptos}
\small
\renewcommand{\arraystretch}{1.4}
\begin{tabular}{p{2.8cm}p{5.2cm}p{7.2cm}}
\hline
\textbf{Concepto} & \textbf{Qué es} & \textbf{Ejemplo} \\
\hline
Hipótesis nula $H_0$ & Afirmación de que no hay diferencia o efecto; se asume cierta hasta que los datos indiquen lo contrario. & $H_0$: la exactitud del modelo nuevo es igual a la del modelo actual. \\
Hipótesis alternativa $H_1$ & Lo que se quiere demostrar; difiere de $H_0$. & $H_1$: la exactitud del modelo nuevo es mayor. \\
Nivel de significancia $\alpha$ & Umbral fijado antes de la prueba; es la probabilidad máxima de error tipo I. & Se fija $\alpha = 0.05$: se tolera rechazar una $H_0$ verdadera en 5 de cada 100 pruebas. \\
Valor $p$ & Compatibilidad de los datos con $H_0$; probabilidad de un resultado tan extremo o más si $H_0$ es cierta. & Tras comparar los modelos se obtiene $p = 0.02$: una diferencia así sería poco común si ambos modelos fueran iguales. \\
Decisión & Se rechaza $H_0$ si $p \le \alpha$; si $p > \alpha$ no se rechaza. & Como $0.02 \le 0.05$, se rechaza $H_0$ y se concluye que el modelo nuevo es mejor. Con $p = 0.20$ no se rechazaría. \\
Error tipo I & Rechazar $H_0$ siendo verdadera. & Concluir que el modelo nuevo es mejor cuando la diferencia fue solo azar. \\
Error tipo II & No rechazar $H_0$ siendo falsa. & No detectar que el modelo nuevo sí es mejor, por ejemplo porque la muestra de prueba es pequeña. \\
\hline
\end{tabular}
\end{table}

\subsection{Decisiones posibles}

Al combinar la realidad (si $H_0$ es verdadera o falsa) con la decisión que toma la prueba resultan cuatro casos, que se muestran en la Tabla~\ref{tab:decisiones}. Solo en dos de ellos la decisión es correcta; en los otros dos se comete uno de los errores del mapa.

\begin{table}[htbp]
\centering
\caption{Resultado de la decisión según la realidad.}
\label{tab:decisiones}
\small
\renewcommand{\arraystretch}{1.5}
\begin{tabular}{p{3.2cm}p{5.8cm}p{5.8cm}}
\hline
 & \textbf{$H_0$ es verdadera} & \textbf{$H_0$ es falsa} \\
\hline
\textbf{Se rechaza $H_0$} & Error tipo I (probabilidad $\alpha$) & Decisión correcta \\
\textbf{No se rechaza $H_0$} & Decisión correcta & Error tipo II (probabilidad $\beta$) \\
\hline
\end{tabular}
\end{table}

\subsection{Explicación de las conexiones}

\textbf{$H_0$ y $H_1$.} Son proposiciones opuestas sobre el mismo parámetro. La prueba parte de suponer que $H_0$ es cierta y busca ver si los datos la contradicen lo bastante como para apoyar $H_1$ \parencite{tecnm2026_documento_fuente}.

\textbf{Valor $p$ y $H_0$.} El valor $p$ se calcula suponiendo que $H_0$ es verdadera, por eso mide la compatibilidad de los datos con ella. Un $p$ cercano a 1 indica datos muy compatibles con $H_0$ y un $p$ cercano a 0 indica datos difíciles de explicar si $H_0$ fuera cierta.

\textbf{Valor $p$ y $\alpha$.} Son los dos lados de la regla de decisión: $p$ es la evidencia que producen los datos y $\alpha$ es el umbral que se fijó antes de verlos. Si $p \le \alpha$ se rechaza $H_0$; si $p > \alpha$ no se rechaza. Conviene decir ``no rechazar'' y no ``aceptar'' $H_0$, porque un $p$ grande solo indica que los datos no bastan para contradecirla.

\textbf{$\alpha$ y el error tipo I.} El nivel de significancia es la probabilidad máxima de cometer el error tipo I, de modo que al elegir $\alpha$ se decide cuánto riesgo de falsa alarma se tolera \parencite{tecnm2026_documento_fuente}. Esa es la razón de que se fije antes de la prueba.

\textbf{Error tipo I y error tipo II.} No se pueden eliminar los dos a la vez. Con una muestra de tamaño fijo, un $\alpha$ más pequeño reduce el error tipo I pero aumenta $\beta$, y viceversa. La única forma de reducir ambos al mismo tiempo es aumentar el tamaño de la muestra \parencite{tecnm2026_documento_fuente}. Cuál error pesa más depende del contexto: si una falsa alarma es muy costosa conviene un $\alpha$ pequeño, y si pasar por alto un efecto real es lo grave conviene proteger contra el error tipo II.

\section{Conclusiones}

Contrastar una hipótesis significa decidir, con datos de una muestra y bajo incertidumbre, entre dos afirmaciones opuestas: $H_0$ (sin diferencia ni efecto) y $H_1$ (lo que se quiere demostrar). La decisión se toma comparando el valor $p$, que indica qué tan compatibles son los datos con $H_0$, con el nivel de significancia $\alpha$, que es el riesgo de error tipo I que se fijó antes de la prueba.

El mapa conceptual muestra que los cinco conceptos forman un solo proceso: las hipótesis definen qué se contrasta, $p$ y $\alpha$ producen la decisión, y los errores tipo I y tipo II describen las dos formas en que esa decisión puede ser incorrecta. Como reducir uno aumenta el otro con una muestra fija, el valor de $\alpha$ debe elegirse según el costo de cada error en el problema, y por eso un resultado ``significativo'' no equivale a una conclusión segura.

\referencias

\end{document}
